\documentclass[11pt]{article}
\usepackage{metricsstyle}
\usepackage{booktabs}
\usepackage{listings}
\lstset{basicstyle=\ttfamily\footnotesize,breaklines=true,columns=fullflexible,showstringspaces=false}

\begin{document}

\hwheader{2 --- Solutions}{Fall 2026}{Xingming Li \quad (Student ID: 1432037)}

% ======================================================================
\question{Question 1}
Data: $X_{1t},X_{2t}\sim U(1,2)$, $u_t\sim N(0,1)$, $n=200$ (seed 6140).
The same $X$'s and $u$'s are used for every value of $\beta$.

\textbf{GNR.} Write the model as $y_t=x_t(\theta)+u_t$ with $x_t(\theta)=\alpha\bigl(X_{1t}^{\beta}+X_{2t}^{1-\beta}\bigr)$,
$\theta=(\alpha,\beta)\xT$. A first-order Taylor expansion around the current value $\theta^{(j)}$ gives
\[
  x_t(\theta)\approx x_t(\theta^{(j)})+X_t(\theta^{(j)})\,(\theta-\theta^{(j)}),
  \qquad X_t(\theta)=\frac{\partial x_t(\theta)}{\partial\theta\xT},
\]
so
\[
  y_t-x_t(\theta^{(j)})=X_t(\theta^{(j)})\,b+\text{resid},\qquad b=\theta-\theta^{(j)} .
\]
This is the GNR. The derivatives are
\begin{align*}
  \frac{\partial x_t}{\partial\alpha}&=X_{1t}^{\beta}+X_{2t}^{1-\beta},\\
  \frac{\partial x_t}{\partial\beta}&=\alpha\Bigl(\frac{\partial X_{1t}^{\beta}}{\partial\beta}+\frac{\partial X_{2t}^{1-\beta}}{\partial\beta}\Bigr)
  =\alpha\bigl(X_{1t}^{\beta}\log X_{1t}-X_{2t}^{1-\beta}\log X_{2t}\bigr),
\end{align*}
using $\partial X^{c}/\partial c=X^{c}\log X$ (and the chain rule gives the minus sign for $1-\beta$).

\emph{Algorithm.} Start from $\theta^{(0)}=(0.5,0.3)$. At step $j$:
\begin{flatnum}
  \item compute the residuals $y_t-x_t(\theta^{(j)})$ and the regressors $X_t(\theta^{(j)})$;
  \item run OLS: $b^{(j)}=\bigl(X\xT X\bigr)^{-1}X\xT\bigl(y-x(\theta^{(j)})\bigr)$;
  \item update $\theta^{(j+1)}=\theta^{(j)}+b^{(j)}$;
  \item stop when $\max|b^{(j)}|<10^{-8}$. At that point $X\xT(y-x(\hat\theta))=0$, which is exactly the
        NLS first-order condition, so the SSR is minimized.
\end{flatnum}
Standard errors come from the final GNR: $\widehat{\Var}(\hat\theta)=s^2(X\xT X)^{-1}$, $s^2=\text{SSR}/(n-2)$.

\begin{center}
\begin{tabular}{llcccc}
\toprule
Part & True $\beta$ & $\hat\alpha$ (s.e.) & $\hat\beta$ (s.e.) & SSR & Iterations \\
\midrule
(i)      & 0.5   & 1.0383 (0.0297) & 0.6126 (0.1680) & 197.1807 & 6 \\
(ii)(a)  & 0.1   & 1.0483 (0.0300) & 0.2279 (0.1633) & 197.1571 & 5 \\
(ii)(b)  & 0.001 & 1.0508 (0.0311) & 0.1305 (0.1607) & 197.1583 & 6 \\
\bottomrule
\end{tabular}
\end{center}

\textbf{Comment.} $\hat\alpha$ is close to its true value 1 and precisely estimated in all three cases.
$\beta$ is estimated imprecisely (s.e.\ $\approx 0.16$), and the smaller the true $\beta$, the further
$\hat\beta$ is pushed away from it (for $\beta=0.001$, $\hat\beta=0.13$). The SSR is nearly identical across
the three regressions, so the objective function is very flat in $\beta$: with $X\in(1,2)$, the
functions $X^{\beta}$ and $X^{1-\beta}$ change very little as $\beta$ changes, so $\beta$ is weakly identified.

% ======================================================================
\question{Question 2}
DGP: $y_t=X_{1t}+X_{2t}+X_{3t}^{2}+u_t$, $X$'s $\sim U(1,2)$, $u_t\sim N(0,1)$.

\textbf{(a)} Here $x_t(\theta)=\beta_1X_{1t}+\beta_2X_{2t}+\beta_3X_{3t}^{\gamma}$, $\theta=(\beta_1,\beta_2,\beta_3,\gamma)\xT$. The derivatives are
\[
  \frac{\partial x_t}{\partial\beta_1}=X_{1t},\qquad
  \frac{\partial x_t}{\partial\beta_2}=X_{2t},\qquad
  \frac{\partial x_t}{\partial\beta_3}=X_{3t}^{\gamma},\qquad
  \frac{\partial x_t}{\partial\gamma}=\beta_3X_{3t}^{\gamma}\log X_{3t},
\]
so each GNR is
\[
  y_t-x_t(\theta^{(j)})=b_1X_{1t}+b_2X_{2t}+b_3X_{3t}^{\gamma^{(j)}}+b_4\,\beta_3^{(j)}X_{3t}^{\gamma^{(j)}}\log X_{3t}+\text{resid}.
\]
Start from $(0.5,0.5,0.5,1.5)$, same updating and stopping rule as Q1, and $s^2=\text{SSR}/(n-4)$.

\begin{center}
\begin{tabular}{lcccccc}
\toprule
 & \multicolumn{2}{c}{$n=50$} & \multicolumn{2}{c}{$n=100$} & \multicolumn{2}{c}{$n=200$} \\
\cmidrule(lr){2-3}\cmidrule(lr){4-5}\cmidrule(lr){6-7}
 & Est. & s.e. & Est. & s.e. & Est. & s.e. \\
\midrule
$\beta_1$ & 0.7125 & 0.4533 & 0.6122 & 0.3431 & 0.8518 & 0.2529 \\
$\beta_2$ & 0.3936 & 0.4807 & 1.3636 & 0.3526 & 0.8542 & 0.2484 \\
$\beta_3$ & 2.4419 & 0.9398 & 1.1101 & 0.6764 & 1.3819 & 0.4637 \\
$\gamma$  & 0.9990 & 0.3332 & 1.8354 & 0.6718 & 1.5958 & 0.3497 \\
\midrule
SSR & \multicolumn{2}{c}{37.9353} & \multicolumn{2}{c}{102.8701} & \multicolumn{2}{c}{212.2670} \\
Iterations & \multicolumn{2}{c}{9} & \multicolumn{2}{c}{9} & \multicolumn{2}{c}{6} \\
\bottomrule
\end{tabular}
\end{center}
Standard errors fall as $n$ grows. $\beta_3$ and $\gamma$ are hard to separate (e.g.\ $n=50$: $\hat\gamma\approx1$
is compensated by $\hat\beta_3=2.44$), since $X_3^{\gamma}$ over $(1,2)$ can be approximated by a rescaled $X_3$.

\textbf{(b)} $H_0:\gamma=1$.

\emph{(i) Unrestricted (Wald / $t$):} $t=(\hat\gamma-1)/\text{se}(\hat\gamma)$, compared with $N(0,1)$.

\emph{(ii) Restricted (LM via GNR):} under $\gamma=1$ the model is linear, so estimate
$y_t=\beta_1X_{1t}+\beta_2X_{2t}+\beta_3X_{3t}+u_t$ by OLS to get $\tilde\beta$, $\tilde u_t$. Then run the GNR
\[
  \tilde u_t = b_1X_{1t}+b_2X_{2t}+b_3X_{3t}+c\,\tilde\beta_3X_{3t}\log X_{3t}+\text{resid},
\]
and use the $t$-statistic on $c$.

\begin{center}
\begin{tabular}{lcccccc}
\toprule
 & \multicolumn{3}{c}{(i) Wald} & \multicolumn{3}{c}{(ii) LM} \\
\cmidrule(lr){2-4}\cmidrule(lr){5-7}
$n$ & $t$ & $p$-value & Reject 5\% & $t_{LM}$ & $p$-value & Reject 5\% \\
\midrule
50  & $-0.0030$ & 0.9976 & No & $-0.0030$ & 0.9976 & No \\
100 & $1.2436$  & 0.2136 & No & $1.8795$  & 0.0602 & No \\
200 & $1.7037$  & 0.0884 & No & $2.3207$  & 0.0203 & Yes \\
\bottomrule
\end{tabular}
\end{center}

\textbf{Comment.} The two tests are asymptotically equivalent but differ in finite samples. At $n=50$ they agree
(since $\hat\gamma\approx1$, the unrestricted and restricted estimates almost coincide) and both fail to reject the
false null. As $n$ grows both statistics increase, but only the LM test rejects $\gamma=1$ at $n=200$; the Wald test
still does not. Both tests have low power here because $\gamma$ is weakly identified.

% ======================================================================
\question{Question 3}

\textbf{(a)} $f(x_i,\theta)=\theta x_i^{\theta-1}$, $0<x_i<1$.

Likelihood and log-likelihood:
\begin{align*}
  \ell(\theta)&=\prod_i\theta x_i^{\theta-1},\\
  L(\theta)&=\sum_i\log\bigl(\theta x_i^{\theta-1}\bigr)=\sum_i\bigl[\log\theta+(\theta-1)\log x_i\bigr]
  =n\log\theta+(\theta-1)\sum_i\log x_i .
\end{align*}
F.O.C.\ w.r.t.\ $\theta$:
\begin{align*}
  \frac{\partial L}{\partial\theta}=\frac{n}{\hat\theta}+\sum_i\log x_i&=0\\
  -\hat\theta\sum_i\log x_i&=n\\
  \hat\theta&=-\frac{n}{\sum_i\log x_i}.
\end{align*}
S.O.C.: $\partial^2L/\partial\theta^2=-n/\theta^2<0$, so this is a maximum. (Since $0<x_i<1$, $\sum_i\log x_i<0$ and $\hat\theta>0$.)

\textbf{(b)} $f(x_i,\theta)=\frac1\theta e^{-x_i/\theta}$, $x_i>0$.
\begin{align*}
  \ell(\theta)&=\prod_i\Bigl(\frac1\theta\Bigr)e^{-x_i/\theta},\\
  L(\theta)&=\sum_i\Bigl[\log\Bigl(\frac1\theta\Bigr)-\frac{x_i}{\theta}\Bigr]
  =n\log\Bigl(\frac1\theta\Bigr)-\frac{\sum_i x_i}{\theta}
  =n\log(1)-n\log\theta-\frac{\sum_i x_i}{\theta}
  =-n\log\theta-\frac{\sum_i x_i}{\theta}.
\end{align*}
F.O.C.\ w.r.t.\ $\theta$:
\begin{align*}
  \frac{\partial L}{\partial\theta}=-\frac{n}{\hat\theta}+\frac{\sum_i x_i}{\hat\theta^2}&=0\\
  \frac1{\hat\theta}\Bigl(-n+\frac{\sum_i x_i}{\hat\theta}\Bigr)&=0\\
  \frac{\sum_i x_i}{\hat\theta}&=n\\
  \hat\theta&=\frac{\sum_i x_i}{n}=\bar x .
\end{align*}
S.O.C.: $\partial^2L/\partial\theta^2\big|_{\hat\theta}=\frac{n}{\hat\theta^2}-\frac{2\sum_ix_i}{\hat\theta^3}=\frac{n}{\hat\theta^2}-\frac{2n}{\hat\theta^2}=-\frac{n}{\hat\theta^2}<0$, a maximum.

\textbf{(c)} $x_i=\theta+u_i^2$, $u_i\sim N(0,\sigma^2)$. Recall the density of $u$:
\[
  f_u(u)=\frac{1}{\sqrt{2\pi\sigma^2}}\exp\Bigl(-\frac{u^2}{2\sigma^2}\Bigr).
\]
\emph{Density of $x$} (Jacobian of the transformation). $x=g(u)=u^2+\theta$ is not one-to-one; just via simple substitution,
\[
  u_i^2=x_i-\theta\quad\Rightarrow\quad u_{1}=+\sqrt{x_i-\theta},\qquad u_{2}=-\sqrt{x_i-\theta},
\]
with
\[
  \frac{du_1}{dx}=\frac12(x_i-\theta)^{-1/2},\qquad \frac{du_2}{dx}=-\frac12(x_i-\theta)^{-1/2}.
\]
Both roots give the same $u^2=x_i-\theta$, so $f_u(u_1)=f_u(u_2)=\frac{1}{\sqrt{2\pi\sigma^2}}\exp\bigl(-\frac{x_i-\theta}{2\sigma^2}\bigr)$. Then
\begin{align*}
  f(x_i)&=f_u(u_1)\Bigl|\frac{du_1}{dx}\Bigr|+f_u(u_2)\Bigl|\frac{du_2}{dx}\Bigr|\\
  &=\frac{1}{\sqrt{2\pi\sigma^2}}\exp\Bigl(-\frac{x_i-\theta}{2\sigma^2}\Bigr)\cdot\frac12(x_i-\theta)^{-1/2}
   +\frac{1}{\sqrt{2\pi\sigma^2}}\exp\Bigl(-\frac{x_i-\theta}{2\sigma^2}\Bigr)\cdot\frac12(x_i-\theta)^{-1/2}\\
  &=\frac{1}{\sqrt{(x_i-\theta)\,2\pi\sigma^2}}\exp\Bigl(-\frac{x_i-\theta}{2\sigma^2}\Bigr),\qquad x_i>\theta .
\end{align*}
\emph{Log-likelihood.}
\begin{align*}
  \ell&=\prod_i f(x_i)=\prod_i\frac{1}{\sqrt{(x_i-\theta)\,2\pi\sigma^2}}\exp\Bigl(-\frac{x_i-\theta}{2\sigma^2}\Bigr),\\
  L&=\sum_i\Bigl[\log(1)-\frac12\log(x_i-\theta)-\frac12\log2\pi-\frac12\log\sigma^2-\frac{x_i-\theta}{2\sigma^2}\Bigr]\\
  &=-\frac12\sum_i\log(x_i-\theta)-\frac n2\log2\pi-\frac n2\log\sigma^2-\frac{\sum_i(x_i-\theta)}{2\sigma^2}.
\end{align*}
F.O.C.\ for $\theta$:
\begin{align*}
  \frac{\partial L}{\partial\theta}=\frac12\sum_i\frac{1}{x_i-\theta}+\frac{n}{2\sigma^2}&=0\\
  \sum_i\frac{1}{x_i-\theta}&=-\frac{n}{\sigma^2}.
\end{align*}
Since $x_i>\theta$ for every $i$, the left-hand side is strictly positive while the right-hand side is negative, so
the FOC has no solution for $\theta$. In fact $\partial L/\partial\theta>0$ everywhere, i.e.\ $L$ is increasing in
$\theta$ over the admissible range $\theta<\min_i x_i$. The ML estimate is therefore the boundary
\[
  \hat\theta=\min_i x_i = x_{(1)} .
\]
F.O.C.\ for $\sigma^2$:
\begin{align*}
  \frac{\partial L}{\partial\sigma^2}=-\frac{n}{2\sigma^2}+\frac{\sum_i(x_i-\theta)}{2\sigma^4}&=0\\
  \frac{1}{2\sigma^4}\sum_i(x_i-\theta)&=\frac{n}{2\sigma^2}\\
  \frac{1}{\sigma^2}&=\frac{n}{\sum_i(x_i-\theta)}\\
  \hat\sigma^2&=\frac{\sum_i(x_i-\hat\theta)}{n}.
\end{align*}

% ======================================================================
\question{Question 4}
Derive the mean and variance of $\hat\sigma^2=e\xT e/n$.

\emph{Step 1: write $e\xT e$ in terms of $u$.} With $M=I_n-X(X\xT X)^{-1}X\xT$ (symmetric, idempotent, $MX=0$),
\begin{align*}
  e&=My=M(X\beta+u)=MX\beta+Mu=Mu,\\
  e\xT e&=u\xT M\xT Mu=u\xT Mu,\\
  \hat\sigma^2&=\frac{u\xT Mu}{n}.
\end{align*}
\emph{Step 2: trace of $M$.} Using $\operatorname{tr}(AB)=\operatorname{tr}(BA)$,
\begin{align*}
  \operatorname{tr}(M)&=\operatorname{tr}(I_n)-\operatorname{tr}\bigl(\underbrace{X}_{n\times k}\underbrace{(X\xT X)^{-1}X\xT}_{k\times n}\bigr)
  =n-\operatorname{tr}\bigl(X\xT X(X\xT X)^{-1}\bigr)\\
  &=n-\operatorname{tr}(I_k)=n-k .
\end{align*}
\emph{Step 3: mean.} $u\xT Mu$ is a scalar, so it equals its trace:
\begin{align*}
  \E(\hat\sigma^2)&=\frac1n\E(u\xT Mu)=\frac1n\E\bigl[\operatorname{tr}(u\xT Mu)\bigr]
  =\frac1n\E\bigl[\operatorname{tr}(Muu\xT)\bigr]\\
  &=\frac1n\operatorname{tr}\bigl(M\,\E(uu\xT)\bigr)=\frac1n\operatorname{tr}(M\cdot\sigma^2I)
  =\frac{\sigma^2}{n}\operatorname{tr}(M)\\
  &=\frac{\sigma^2}{n}(n-k)=\sigma^2-\frac{\sigma^2k}{n}=\sigma^2\Bigl(1-\frac kn\Bigr).
\end{align*}
\emph{Step 4: variance.} Define
\[
  w=\frac{e\xT e}{\sigma^2}=\Bigl(\frac{u}{\sigma}\Bigr)\xT M\Bigl(\frac{u}{\sigma}\Bigr)\sim\chi^2_{n-k},
\]
since $u/\sigma\sim N(0,I_n)$ and $M$ is idempotent with $\rank(M)=\operatorname{tr}(M)=n-k$. Then
\[
  \hat\sigma^2=\frac{\sigma^2}{n}\cdot\frac{e\xT e}{\sigma^2}=\frac{\sigma^2}{n}\,w ,
\]
and using $\E(w)=n-k$, $\Var(w)=2(n-k)$:
\begin{align*}
  \E(\hat\sigma^2)&=\frac{\sigma^2}{n}\E(w)=\frac{\sigma^2}{n}(n-k),\\
  \Var(\hat\sigma^2)&=\Bigl(\frac{\sigma^2}{n}\Bigr)^2\Var(w)=\frac{\sigma^4}{n^2}\cdot2(n-k)=\frac{2\sigma^4(n-k)}{n^2}.
\end{align*}
So $\hat\sigma^2$ is biased downward, but the bias vanishes as $n\to\infty$.

% ======================================================================
\question{Question 5}
\begin{align*}
  \ell(\alpha)&=\prod_i\frac{\alpha\,10000^{\alpha}}{y_i^{\alpha+1}}=\prod_i\alpha\,10000^{\alpha}y_i^{-\alpha-1},\\
  L(\alpha)&=\sum_i\bigl[\log\alpha+\alpha\log10000+(-\alpha-1)\log y_i\bigr]\\
  &=n\log\alpha+n\alpha\log10000-(\alpha+1)\sum_i\log y_i .
\end{align*}
F.O.C.\ w.r.t.\ $\alpha$:
\begin{align*}
  \frac{\partial L}{\partial\alpha}=\frac{n}{\hat\alpha}+n\log10000-\sum_i\log y_i&=0\\
  \frac{n}{\hat\alpha}&=\sum_i\log y_i-n\log10000\\
  \hat\alpha&=\frac{n}{\sum_i\log y_i-n\log10000}=\frac{n}{\sum_i\log(y_i/10000)} .
\end{align*}
S.O.C.: $\partial^2L/\partial\alpha^2=-n/\alpha^2<0$, a maximum. (Since $y_i\ge10000$, the denominator is positive.)

% ======================================================================
\question{Question 6}
\textbf{(i)} Let $r_t=y_t-X_t\beta$, so $u_t=r_t^{\lambda}$.

\emph{Log-likelihood.} The density of $u_t$ is $\frac{1}{\sqrt{2\pi\sigma^2}}\exp\bigl(-\frac{u_t^2}{2\sigma^2}\bigr)$.
To get the density of $y_t$ we need the Jacobian
\[
  \Bigl|\frac{\partial u_t}{\partial y_t}\Bigr|=\bigl|\lambda r_t^{\lambda-1}\bigr|=\lambda|r_t|^{\lambda-1},
\]
so
\[
  f(y_t)=\frac{1}{\sqrt{2\pi\sigma^2}}\exp\Bigl(-\frac{|r_t|^{2\lambda}}{2\sigma^2}\Bigr)\lambda|r_t|^{\lambda-1},
\]
and summing the logs over $t$,
\[
  L=-\frac n2\log2\pi-\frac n2\log\sigma^2-\frac1{2\sigma^2}\sum_t|r_t|^{2\lambda}+n\log\lambda+(\lambda-1)\sum_t\log|r_t| .
\]
\emph{F.O.C.\ for $\beta$.} Since $\partial r_t/\partial\beta=-X_t\xT$, for safety write $|r_t|^{2\lambda}=(r_t^2)^{\lambda}$ and use the chain rule:
\begin{align*}
  \frac{\partial}{\partial\beta}\Bigl[\frac1{2\sigma^2}\sum_t(r_t^2)^{\lambda}\Bigr]
  &=\frac1{2\sigma^2}\sum_t\frac{\partial(r_t^2)^{\lambda}}{\partial r_t^2}\cdot\frac{\partial r_t^2}{\partial r_t}\cdot\frac{\partial r_t}{\partial\beta}\\
  &=\frac1{2\sigma^2}\sum_t\lambda(r_t^2)^{\lambda-1}\cdot2r_t\cdot(-X_t\xT)
  =-\frac{\lambda}{\sigma^2}\sum_t(r_t^2)^{\lambda-1}r_tX_t\xT,
\end{align*}
and
\[
  \frac{\partial}{\partial\beta}\Bigl[(\lambda-1)\sum_t\log|r_t|\Bigr]=(\lambda-1)\sum_t\frac{1}{r_t}\cdot(-X_t\xT)=-(\lambda-1)\sum_t\frac{X_t\xT}{r_t}.
\]
Hence
\[
  \frac{\partial L}{\partial\beta}=\frac{\lambda}{\sigma^2}\sum_t(r_t^2)^{\lambda-1}r_tX_t\xT-(\lambda-1)\sum_t\frac{X_t\xT}{r_t}=0 ,
\]
which has no closed-form solution for $\beta$.

\emph{F.O.C.\ for $\sigma^2$:}
\begin{align*}
  \frac{\partial L}{\partial\sigma^2}=-\frac{n}{2\sigma^2}+\frac1{2\sigma^4}\sum_t|r_t|^{2\lambda}&=0\\
  \frac{n}{2\sigma^2}&=\frac{\sum_t|r_t|^{2\lambda}}{2\sigma^4}\\
  n\sigma^2&=\sum_t|y_t-X_t\beta|^{2\lambda}\\
  \hat\sigma^2&=\frac{\sum_t|y_t-X_t\hat\beta|^{2\lambda}}{n}.
\end{align*}
\emph{If $\lambda$ is known}, solve these two FOCs for $\hat\beta,\hat\sigma^2$. There is no explicit solution for $\beta$,
so it is done numerically (e.g.\ substitute $\hat\sigma^2(\beta)$ back into $L$ and maximize the concentrated
log-likelihood over $\beta$). For $\lambda=1$ the Jacobian term drops out and this is just OLS.

\emph{If $\lambda$ is unknown}, we also need the F.O.C.\ for $\lambda$. Using $\partial|r_t|^{2\lambda}/\partial\lambda=2|r_t|^{2\lambda}\log|r_t|$,
\begin{align*}
  \frac{\partial L}{\partial\lambda}
  &=\frac{n}{\lambda}+\sum_t\log|r_t|-\frac1{2\sigma^2}\sum_t2\,|r_t|^{2\lambda}\log|r_t|\\
  &=\frac{n}{\lambda}+\sum_t\log|r_t|-\frac1{\sigma^2}\sum_t|r_t|^{2\lambda}\log|r_t|=0 .
\end{align*}
Solve the three FOCs jointly to get $\hat\beta$, $\hat\sigma^2$ and $\hat\lambda$, numerically on the computer.

\medskip
\textbf{(ii)} $H_0:\lambda=1$, using only the restricted estimates (LM test via the OPG regression).

\emph{Restricted estimates.} With $\lambda=1$, $n\log\lambda=0$, $(\lambda-1)\sum_t\log|r_t|=0$ and $|r_t|^{2\lambda}=r_t^2$, so
\[
  L(\beta,\sigma^2,1)=-\frac n2\log2\pi-\frac n2\log\sigma^2-\frac1{2\sigma^2}\sum_t(y_t-X_t\beta)^2 ,
\]
which is the normal linear regression model. Hence
\[
  \tilde\beta=(X\xT X)^{-1}X\xT y,\qquad e_t=y_t-X_t\tilde\beta,\qquad \tilde\sigma^2=\frac{e\xT e}{n}.
\]
\emph{Gradient contributions.} Write $L=\sum_t\ell_t$ with
$\ell_t=-\frac12\log2\pi-\frac12\log\sigma^2-\frac{(r_t^2)^\lambda}{2\sigma^2}+\log\lambda+(\lambda-1)\log|r_t|$.
Differentiating as in (i) and evaluating at $\tilde\theta$ ($\lambda=1$, $r_t=e_t$):
\begin{align*}
  \frac{\partial\ell_t}{\partial\beta}&=\frac{\lambda}{\sigma^2}(r_t^2)^{\lambda-1}r_tX_t\xT-(\lambda-1)\frac{X_t\xT}{r_t}
  &&\Rightarrow\quad \frac{e_tX_t\xT}{\tilde\sigma^2},\\
  \frac{\partial\ell_t}{\partial\sigma^2}&=-\frac1{2\sigma^2}+\frac{(r_t^2)^\lambda}{2\sigma^4}
  &&\Rightarrow\quad -\frac1{2\tilde\sigma^2}+\frac{e_t^2}{2\tilde\sigma^4},\\
  \frac{\partial\ell_t}{\partial\lambda}&=\frac1\lambda+\log|r_t|-\frac{(r_t^2)^\lambda}{\sigma^2}\log|r_t|
  &&\Rightarrow\quad 1+\log|e_t|-\frac{e_t^2}{\tilde\sigma^2}\log|e_t| ,
\end{align*}
using $\partial(r_t^2)^\lambda/\partial\lambda=2(r_t^2)^\lambda\log|r_t|$. Stack these into the $n\times(k+2)$ matrix $\tilde G$ (row $t$ = the three
contributions above).

\emph{Score.} $\tilde g=\tilde G\xT\iota$. The $\beta$ block is $X\xT e/\tilde\sigma^2=0$ (normal equations) and the $\sigma^2$ element is
$-n/(2\tilde\sigma^2)+e\xT e/(2\tilde\sigma^4)=0$, so only
\[
  \tilde g_\lambda=n+\sum_t\log|e_t|-\frac1{\tilde\sigma^2}\sum_te_t^2\log|e_t|
\]
is nonzero. The test asks whether $\tilde g_\lambda$ is significantly different from 0.

\emph{LM statistic.} Estimating the information matrix by $\tilde G\xT\tilde G$ (OPG),
\[
  LM=\tilde g\xT(\tilde G\xT\tilde G)^{-1}\tilde g=\iota\xT\tilde G(\tilde G\xT\tilde G)^{-1}\tilde G\xT\iota=\text{ESS}=n-\text{SSR},
\]
from the regression of $\iota$ (a column of ones) on all $k+2$ columns of $\tilde G$, no constant.
All columns must be included, so that the variance of $\tilde g_\lambda$ accounts for $\beta,\sigma^2$ being estimated.
Under $H_0$, $LM\overset{a}{\sim}\chi^2(1)$; reject at 5\% if $LM>3.84$.

\emph{Procedure.} (1) OLS of $y$ on $X$, save $e_t$, $\tilde\sigma^2$; (2) build the three columns of $\tilde G$;
(3) regress $\iota$ on $\tilde G$; (4) $LM=n-\text{SSR}$ vs.\ $\chi^2_{0.95}(1)=3.84$. The model with $\lambda$ free is never estimated.

% ======================================================================
\newpage
\topic{Appendix: Python code (Questions 1--2)}
\lstinputlisting[language=Python]{hw2.py}

\end{document}
